% source: David Mumford, "Abelian Varieties", TIFR Studies in Mathematics 5, Oxford UP, 1970
% pdf_page: 0055
% doc_page: 44 (Chapter II, Section 4, corollaries of the rigidity lemma; Appendix to Section 4)
% retrieved: 2026-07-11
% transcribed_by: claude-fable-5 (vision, rendered at 200dpi via PyMuPDF; scanned book, no text layer)

\DeclareMathOperator{\Hom}{Hom}

% [running head: ABELIAN VARIETIES, page 44]

Note that in the proof of the above corollary, despite the additive notation
used, no use of the commutativity of $X$ was made. Thus, we may use it to
give a second proof of the commutativity of $X$.

\paragraph{Corollary 2.}
\textit{$X$ is a commutative group.}

In fact, the morphism of $X$ into itself mapping each element onto its
inverse is a homomorphism by Corollary 1. Hence, for $x, y \in X$,
$(xy)^{-1} = x^{-1} y^{-1} = y^{-1} x^{-1}$, and $X$ is commutative.

\paragraph{Corollary 3.}
\textit{Let $X$ be an abelian variety (with base point $0$). Then on the
category of complete varieties with base point, the functor
$S \longmapsto \Hom(S, X)$ (where $\Hom$ denotes morphisms preserving base
point) is linear; that is, for $S$, $T$ in this category, the natural map}
\[
  \Hom(S, X) \times \Hom(T, X) \longrightarrow \Hom(S \times T, X)
\]
\textit{given by $(f, g) \longmapsto h$, $h(s,t) = f(s) + g(t)$ is a
bijection.}

\paragraph{Proof.}
That this map is injective follows by fixing $s$ and $t$ in turn to be the
base points $s_0$ and $t_0$ of $S$ and $T$, respectively, in the equation
$h(s, t) = f(s) + g(t)$. Next, take any $h \in \Hom(S \times T, X)$, and put
$f(s) = h(s, t_0)$, $g(t) = h(s_0, t)$,
$k(s, t) = h(s, t) - f(s) - g(t)$. Then
$k(S \times \{t_0\}) = k(\{s_0\} \times T) = 0$, and it follows from the
rigidity lemma that $k \equiv 0$.

\section*{Appendix to \S 4}

We want to prove the following result.

\paragraph{Theorem.}
\textit{Let $X$ be a complete variety, $e \in X$ a point, and}
\[
  m : X \times X \longrightarrow X
\]
\textit{a morphism such that $m(x, e) = m(e, x) = x$ for all $x \in X$. Then
$X$ is an abelian variety with group law $m$ and identity $e$.}

\paragraph{Proof.}
We shall denote $m(x, y)$ simply by $xy$. Introduce the morphism
\[
  \psi : X \times X \longrightarrow X \times X
\]
\[
  \psi : (x, y) = (xy, y).
\]
% [proof continues on the next page]
